\section{Metodologia}
\label{sec:metodologia}

A metodologia seguiu diretamente a formulação dos mínimos quadrados apresentada
nas Seções iniciais. As bases polinomiais descritas em $\S$\ref{sec:regressao}
foram geradas tanto pela abordagem simbólica quanto pela implementação
vetorizada do pacote \texttt{regkit}, garantindo que os sistemas normais
$G\vec{a}=\vec{b}$ produzidos em cada experimento correspondessem à teoria. A
versão simbólica atuou como referência para depuração, enquanto a variante
otimizada foi empregada nas execuções de larga escala descritas a seguir.

\subsection{Problemas estudados}

Foram analisados três problemas complementares. O primeiro é o conjunto
exploratório de hélices gerado por \texttt{gen\_exploration\_dataset}, com
$n=3$, $m=2$, $N=200$ e ruído gaussiano de desvio padrão $0{,}01$. Ele codifica
trajetórias helicoidais acopladas entre entradas e saídas e serve como cenário
controlado para validar a montagem matricial e comparar as implementações
ingênua e vetorizada. O segundo problema corresponde ao protótipo~4$\times$4,
conjunto alvo sintético obtido por \texttt{create\_target\_dataset} que
reproduz a versão estruturada em MATLAB com quatro entradas e quatro saídas
parametrizadas pelos eixos temporais $t$ e $q$. A heterogeneidade entre os
regimes de tempo e parâmetro torna esse protótipo o foco principal da aplicação.
O problema extra é o ajuste da superfície ``wing'' a partir de 308 amostras
avaliadas sobre a geometria da asa. 

\subsection{Procedimentos experimentais e visualização}

Cada execução foi gerenciada pelos scripts \texttt{run\_regkit\_usage.py} e
\texttt{run\_benchmark\_regression.py}. O primeiro instancia os problemas das
hélices e do protótipo~4$\times$4, executa um \textit{sweep} nos graus polinomiais de
1 a~10 e persiste as métricas de erro em \texttt{results/regkit\_usage\_rmse.csv}.
O segundo consolida os experimentos da asa e os benchmarks de desempenho. A
análise qualitativa utiliza o utilitário \texttt{plot\_dataset}, que aplica PCA
em $\mathbb{R}^{n+m}$ e reaproveita o mesmo projetor para comparar dados
originais e predições. Esse procedimento mantém a
correspondência geométrica entre séries reais e estimadas, além de orientar a
leitura dos resíduos nas figuras apresentadas na Seção~\ref{sec:resultados}. 
Vale ressaltar que a utilização do PCA é necessária para visualizar a semelhança 
entre os conjuntos, dada a alta dimensionalidade dos dados. Apenas o problema da asa
permite uma visualização direta, pois ocorre em três dimensões. 
